\chapter{Qualité des données et lignage}\label{ch:pl:data-quality-and-lineage}

Le 6 mai 2010, entre 14~h~40 et 15~h~00, plus de 20\,000 transactions portant sur plus de 300 titres américains ont été exécutées à
des cours s'écartant de 60~\% ou plus de ceux qui avaient été affichés quelques minutes plus tôt, certaines à un cent, d'autres à 100\,000~\$. Après la clôture, les
bourses et leur autorégulateur ont convenu d'annuler toutes ces transactions comme manifestement erronées. Une \omterm{def:pl:a-tick-store-on-open-formats:store}{base de ticks} qui avait enregistré ces
transactions fidèlement, sans jamais apprendre qu'elles avaient été annulées, les a transmises à chaque backtest, à chaque estimation
de volatilité et à chaque étude d'exécution portant sur cet après-midi, tant qu'elle a conservé ces données. Rien dans la base n'était corrompu~; les données étaient
exactement ce que le flux avait transmis. Il manquait une couche qui confronte les données à ce qu'elles devraient être, signale ce qui
échoue au lieu de la supprimer, applique les corrections qui arrivent plus tard sous forme de nouvelles versions, et sait quels résultats ont été
construits à partir de quoi. Ce chapitre construit cette couche (\texttt{firm.dataqual}) et la mesure sur une journée simulée comportant des défauts introduits et,
tout aussi important, de véritables événements de marché qu'une règle peu rigoureuse prend pour des défauts.

\section{Contrôles~: complétude, validité, cohérence, ponctualité}

\begin{definition}[Règle de qualité des données, contrat de données]\label{def:pl:data-quality-and-lineage:rule}
Une \emph{règle de qualité des données}\index{règle de qualité des données} est un contrôle automatisé d'un jeu de données par rapport à une attente~--- qu'il est
complet, que ses valeurs sont valides, qu'il est cohérent avec lui-même et avec d'autres sources, qu'il est ponctuel~--- exécutée
à chaque fois que les données sont produites, avec une gravité et un responsable qui agit sur ses échecs. Un \emph{contrat de données}\index{contrat de données}
est l'accord entre le producteur d'un jeu de données et ses consommateurs, qui en fixe le schéma, les règles et les délais de livraison, afin que
une modification de l'un ou l'autre côté soit traitée comme une modification du contrat plutôt que comme une surprise.
\end{definition}

Quatre familles de contrôles couvrent l'essentiel des défauts des données de marché. La complétude~: aucun message ne manque (les numéros de séquence du flux ne présentent aucun
trou, chapitre~2), aucun instrument silencieux alors qu'il devrait être coté. Validité~: chaque valeur reste dans son domaine (prix et tailles positifs,
instruments connus, horodatages compris dans la séance). Cohérence~: les données concordent entre elles (aucun cours acheteur égal ou supérieur au cours vendeur, transactions
à des prix proches des cotations) et avec d'autres sources (un autre flux, la copie d'un fournisseur, le cours de clôture officiel). Ponctualité~: les données sont arrivées
au moment prévu par le contrat. Le \Cref{lst:pl:data-quality-and-lineage:rules} est la suite de règles du chapitre, écrite sous forme de données~: chaque règle a
un type, une gravité, un responsable et des paramètres, et \texttt{firm.dataqual} exécute n'importe quelle suite sur n'importe quelle partition.

\omcode{code/platforms/07-data-quality-and-lineage/python/pl_dataqual.py}{88}{110}{La suite de règles du chapitre, sous forme de données~: trous de séquence,
bornes, carnets croisés, cotations périmées, pics de prix, transactions annulées et un décompte par rapport à la copie d'un fournisseur, chacun avec une gravité et un
responsable.}\label{lst:pl:data-quality-and-lineage:rules}

La journée est la deuxième de la semaine du chapitre~4~: deux instruments, 47\,472 événements, 9\,258 cotations et 1\,190 transactions. Deux
événements réels en font partie~: SIM2 fait l'objet d'une suspension de cotation de 60 secondes à 40~\% de la séance, et le niveau de prix de SIM2 bondit de 3~\% à 60~\% de celle-ci,
à la suite d'une annonce, et s'y maintient. Dans des copies de la journée, vingt fois avec vingt graines aléatoires, le chapitre introduit six types de défauts~: un trou de
200 événements, trois cotations croisées, huit prix de transaction décalés d'un facteur dix (une erreur de virgule, à la hausse ou à la baisse), un silence de 90 secondes des cotations de
SIM1, quatre transactions annulées par la plateforme après la clôture et, séparément, une \omterm{def:pl:data-quality-and-lineage:correction}{correction du fournisseur} portant sur un cours de clôture officiel.

\section{Indicateurs d'anomalie plutôt que suppressions}

\begin{definition}[Indicateur d'anomalie, quarantaine]\label{def:pl:data-quality-and-lineage:flag}
Un \emph{indicateur d'anomalie}\index{indicateur d'anomalie} est un enregistrement rattaché à une ligne ou à une partition, qui désigne la règle en échec, la
gravité et le détail, sans modifier la ligne elle-même~; les consommateurs décident s'ils excluent les lignes signalées. Une
\emph{quarantaine}\index{quarantaine} est l'état d'une partition ayant échoué à une règle de gravité «~erreur~»~: elle est tenue
à l'écart des lecteurs jusqu'à ce qu'une personne la libère, avec un motif, ou qu'elle soit reconstruite.
\end{definition}

Supprimer une mauvaise ligne est tentant et erroné. La ligne peut être juste (un saut réel de 3~\% ressemble à un pic pour une règle naïve)~; un autre
consommateur peut en avoir besoin (un système de surveillance veut les transactions erronées, chapitre~30)~; et une suppression ne peut être ni auditée ni annulée.
Un indicateur garde les trois options ouvertes~: le backtest exclut les lignes de gravité «~erreur~», le système de surveillance les lit et la règle
peut être améliorée et relancée.

\begin{omfigure}
\begin{tikzpicture}[font=\scriptsize,>=Stealth,box/.style={draw,rounded corners=2pt,fill=omDef!10,minimum height=0.7cm,inner sep=3pt,align=center}]
\node[box] (p) at (0,0) {partition\\(date, table)};
\node[box,fill=omMeth!15] (r) at (2.9,0) {suite de règles\\(données)};
\node[box,fill=omThm!12] (f) at (5.7,0) {indicateurs\\(ligne, règle,\\gravité)};
\node[box] (pub) at (9.0,1.0) {publié~:\\les consommateurs\\filtrent};
\node[box,fill=omThm!20] (qz) at (9.0,-1.0) {en quarantaine~:\\retenue};
\node[box,fill=gray!10] (rel) at (12.1,-1.0) {libérée\\par une personne};
\node[box,fill=omProp!15] (cor) at (2.9,-1.8) {correction du fournisseur~:\\nouvelle version};
\draw[->] (p) -- (r); \draw[->] (r) -- (f);
\draw[->] (f) -- node[above,sloped,font=\tiny]{avertissements} (pub.west);
\draw[->] (f) -- node[below,sloped,font=\tiny]{erreurs} (qz.west);
\draw[->] (qz) -- (rel);
\draw[->] (cor) -| (p);
\end{tikzpicture}

\omcaption{La couche de qualité des données~: chaque partition passe la suite de règles~; les indicateurs sont enregistrés à côté des lignes, jamais à leur place~;
une partition portant un indicateur de gravité «~erreur~» est mise en \omterm{def:pl:data-quality-and-lineage:flag}{quarantaine} jusqu'à sa libération~; une \omterm{def:pl:data-quality-and-lineage:correction}{correction du fournisseur} arrive sous forme de nouvelle version des données,
et les règles sont à nouveau exécutées.}\label{fig:pl:data-quality-and-lineage:pipeline}
\end{omfigure}

\begin{omfigure}
\begin{tikzpicture}
\begin{axis}[width=12cm,height=5.6cm,ymode=log,xlabel={numéro de transaction de SIM2, dans l'ordre de séquence},ylabel={prix de transaction (échelle logarithmique)},
  tick label style={font=\small},label style={font=\small},grid=major,grid style={gray!20},ymin=20,ymax=5000,
  ytick={30,100,300,1000,3000},yticklabels={30,100,300,1\,000,3\,000},
  legend columns=2,legend style={font=\footnotesize,at={(0.5,-0.3)},anchor=north,draw=none,/tikz/every even column/.append style={column sep=8pt}},
  table/col sep=comma]
\addplot[only marks,mark=*,mark size=0.8pt,omDef] table[x=n,y=price]{figdata/platforms/07-data-quality-and-lineage/sim2_prices.csv};
\addplot[only marks,mark=o,mark size=3pt,omThm,thick] table[x=n,y=price,restrict expr to domain={\thisrow{flagged}}{1:1}]{figdata/platforms/07-data-quality-and-lineage/sim2_prices.csv};
\legend{transaction,signalée comme pic (erreur)}
\end{axis}
\end{tikzpicture}

\omcaption{Les transactions de SIM2 sur la copie défectueuse de la graine~0, en échelle logarithmique~: deux erreurs de virgule (l'une dix fois trop haute, l'autre dix fois trop
basse) ressortent et sont signalées~; le saut réel de 3~\% survenu en fin de journée est invisible à cette échelle et n'est pas signalé par la règle qui regarde aussi vers l'avant.
Données~: \texttt{fig\_dataqual.py}.}\label{fig:pl:data-quality-and-lineage:prices}
\end{omfigure}

La première version de la suite a signalé les défauts introduits et bien davantage. La règle des pics comparait chaque prix de transaction avec la
médiane des 50 transactions précédentes de son symbole, en unités robustes~: l'écart absolu médian des écarts précédents, avec un
plancher de 0.1~\% afin qu'un mouvement d'un seul pas de cotation ne soit jamais un pic. Le jour du saut, elle a signalé chacune des quelques dizaines de transactions suivantes, jusqu'à ce que la médiane rétrospective rattrape le
nouveau niveau (\cref{fig:pl:data-quality-and-lineage:jump}). Un pic et un changement de niveau se ressemblent vus de l'arrière~; vus de l'avant, ils
ne se ressemblent plus, car un pic revient. L'exécution des contrôles de qualité a lieu après la clôture, elle peut donc regarder vers l'avant~:

\begin{method}[Distinguer un pic d'un mouvement]\label{met:pl:data-quality-and-lineage:spike}
Signaler une valeur lorsqu'elle s'écarte, de plus de $k$ unités robustes, à la fois de la médiane des valeurs qui la précèdent et de la médiane de
les valeurs qui la suivent. Une erreur de virgule s'écarte des deux~; un mouvement réel du niveau ne s'écarte que des valeurs qui précèdent. Au
début d'une série, où il n'y a pas d'historique, on ne compare qu'aux valeurs suivantes.
\end{method}

\omcode{code/firm/dataqual/firm_dataqual.py}{78}{99}{La règle des pics~: écart à la médiane rétrospective en unités robustes, avec plancher,
et, avec \texttt{confirm}, par rapport à la médiane prospective également.}\label{lst:pl:data-quality-and-lineage:spike}

\begin{omfigure}
\begin{tikzpicture}
\begin{axis}[width=12cm,height=5.8cm,xlabel={numéro de transaction de SIM2},ylabel={prix},tick label style={font=\small},
  label style={font=\small},grid=major,grid style={gray!20},enlarge y limits=0.12,
  legend columns=2,legend style={font=\footnotesize,at={(0.5,-0.3)},anchor=north,draw=none,/tikz/every even column/.append style={column sep=8pt}},
  table/col sep=comma]
\addplot[only marks,mark=*,mark size=1.2pt,omDef] table[x=n,y=price]{figdata/platforms/07-data-quality-and-lineage/jump.csv};
\addplot[thick,omMeth] table[x=n,y=behind]{figdata/platforms/07-data-quality-and-lineage/jump.csv};
\addplot[thick,omMeth,dashed] table[x=n,y=ahead]{figdata/platforms/07-data-quality-and-lineage/jump.csv};
\addplot[only marks,mark=o,mark size=2.6pt,omThm,thick] table[x=n,y=price,restrict expr to domain={\thisrow{naive_only}}{1:1}]{figdata/platforms/07-data-quality-and-lineage/jump.csv};
\legend{prix de transaction,médiane des 50 précédentes,médiane des 50 suivantes,signalée par la seule règle naïve}
\end{axis}
\end{tikzpicture}

\omcaption{Les transactions de SIM2 autour de son véritable saut de 3~\%. Mesurée à partir de la médiane rétrospective, chaque transaction des quelques dizaines suivantes
ressemble à un pic, et la règle naïve en signale 24 (entourées)~; mesurée aussi à partir de la médiane prospective, aucune ne l'est, puisque le nouveau niveau persiste.
Données~: \texttt{fig\_dataqual.py} (graine 0).}\label{fig:pl:data-quality-and-lineage:jump}
\end{omfigure}

La règle de fraîcheur échoue de la même façon sur la suspension de cotation~: 60 secondes sans cotation de SIM2 est un défaut du flux que si le marché
était ouvert. La règle doit connaître l'état du marché --- les messages d'état de la négociation émis par la plateforme (les événements du chapitre~2 les incluent) --- et
ne pas compter le temps passé en suspension de cotation (le \cref{lst:pl:data-quality-and-lineage:stale}).

\omcode{code/firm/dataqual/firm_dataqual.py}{102}{117}{La règle de fraîcheur, qui ne compte pas le temps que le marché lui-même a
passé en silence (suspension de cotation, pause programmée) dans un intervalle.}\label{lst:pl:data-quality-and-lineage:stale}

Le \Cref{tab:pl:data-quality-and-lineage:detection} donne le résultat sur vingt graines. Chaque défaut injecté est trouvé par les deux versions de
la batterie de tests~; la différence tient entièrement aux faux signalements~: 495 transactions signalées à tort comme pics et 20 suspensions de cotation signalées à tort comme
flux périmés par la batterie naïve, aucun par la batterie consciente.

\begin{table}[ht]
\centering\small
\begin{tabular}{lrrrr}
\toprule
règle & injectés & trouvés & faux signalements, naïve & faux signalements, consciente \\
\midrule
rupture de séquence & 20 & 20 & 0 & 0 \\
carnet croisé & 60 & 60 & 0 & 0 \\
pic de prix & 160 & 160 & 495 & 0 \\
cotations périmées & 20 & 20 & 20 & 0 \\
transactions invalidées & 80 & 80 & 0 & 0 \\
\bottomrule
\end{tabular}
\caption{Défauts injectés, défauts trouvés et lignes signalées à tort par chaque règle sur vingt graines d'injection, avec la batterie naïve
(règle de pic rétrospective, règle de fraîcheur aveugle aux suspensions) et la batterie consciente. Données~: \texttt{pl\_dataqual.detection}.}\label{tab:pl:data-quality-and-lineage:detection}
\end{table}

Un défaut injecté que toutes les règles trouvent ne teste pas les règles~; ce qui les éprouve, ce sont les facteurs de confusion. Un détecteur de
surveillance du livre~9 a repéré tous ses spoofeurs injectés tant que des imitations légitimes n'ont pas été ajoutées (son chapitre~29), et la même leçon vaut ici~:
il faut injecter les événements authentiques que comporte le marché réel, et mesurer les faux signalements.

\section{Contrôles inter-sources}

\begin{definition}[Contrôle inter-sources]\label{def:pl:data-quality-and-lineage:cross}
Un \emph{contrôle inter-sources}\index{contrôle inter-sources} compare un jeu de données à une source indépendante des mêmes faits ---
l'autre ligne de flux, un second fournisseur, le fichier de fin de journée de la plateforme, le cours de clôture officiel --- sur un agrégat que les deux peuvent calculer (un nombre, un
volume, un dernier cours par intervalle), et signale les intervalles où elles diffèrent de plus d'une tolérance.
\end{definition}

Une source unique peut être complète selon ses propres numéros de séquence et rester fausse~: un normaliseur qui a écarté un type de message, un
compactage qui a perdu un fichier. Le contrôle du chapitre compte les transactions par minute dans la base et dans une copie du jour fournie par un fournisseur, et signale
les minutes qui diffèrent de plus de 2~\%~: sur les vingt graines, il a levé 16 signalements, dans les minutes où le trou injecté de 200 événements
avait par hasard contenu des transactions. Il ne dit pas quelles transactions manquent, seulement où regarder~; c'est la bonne répartition du travail
entre un contrôle peu coûteux qui tourne partout et une enquête coûteuse qui se déclenche là où il pointe.

\section{Corrections du fournisseur et versions}

\begin{definition}[Correction du fournisseur]\label{def:pl:data-quality-and-lineage:correction}
Une \emph{correction du fournisseur}\index{correction du fournisseur} est une livraison ultérieure qui modifie une valeur déjà livrée (un cours de clôture, une
condition de transaction, une opération sur le capital)~: elle est stockée comme une nouvelle version du même fait, avec son propre temps de connaissance, et n'écrase jamais
la version antérieure.
\end{definition}

Le fournisseur du chapitre livre le cours de clôture officiel de SIM2 à 16~h~30 le jour même, et un cours de clôture corrigé à 7~h~10 le lendemain matin. Les deux sont
des versions d'un même fait dans la base bitemporelle du livre~7~: interrogée à 23~h, la clôture est la première~; interrogée à 8~h, la seconde. Les
transactions invalidées de l'accroche du chapitre suivent le même schéma~: une liste d'annulations, livrée après la clôture, appliquée sous forme d'indicateurs (\texttt{busted}) sur les
transactions qu'elles désignent. Ce qu'une correction modifie n'est pas seulement le fait, mais tout ce qui en est calculé, et c'est précisément le lignage.

\section{Lignage et provenance}

\begin{definition}[Lignage des données, fiche de provenance]\label{def:pl:data-quality-and-lineage:lineage}
Le \emph{lignage des données}\index{lignage des données} d'un résultat est le graphe des versions de jeux de données et des transformations dont il est issu, en remontant
jusqu'à ses entrées brutes. Une \emph{fiche de provenance}\index{fiche de provenance} est la description enregistrée de la manière dont une version de jeu de données a été
produite~: les versions de ses entrées (par empreinte de contenu), le code et sa version, ainsi que le moment de son exécution et son auteur.
\end{definition}

Le lignage de \texttt{firm.dataqual} enregistre chaque version de jeu de données comme un nœud, nommé par l'empreinte de contenu de ses données
(calculée par \texttt{firm.workflow.content\_hash} du livre~7), avec une arête depuis chaque version d'entrée dont
elle a été construite (le \cref{lst:pl:data-quality-and-lineage:lineage}). Une correction produit une nouvelle version de l'entrée~; tout nœud situé en aval de
l'ancienne version devient périmé. Dans l'exemple du chapitre (la \cref{fig:pl:data-quality-and-lineage:lineage}), la clôture corrigée rend
périmées les valorisations, le P\&L et les variables de recherche, et laisse intactes les barres à la minute, construites à partir des seules transactions.

\omcode{code/firm/dataqual/firm_dataqual.py}{189}{203}{Lignage~: une version de jeu de données est son nom et son empreinte de contenu, avec les versions de ses entrées
versions~; tout ce qui se trouve en aval d'une ancienne version devient périmé une fois celle-ci corrigée.}\label{lst:pl:data-quality-and-lineage:lineage}

\begin{omfigure}
\begin{tikzpicture}[font=\scriptsize,>=Stealth,box/.style={draw,rounded corners=2pt,fill=omDef!10,minimum height=0.6cm,inner sep=3pt,align=center}]
\node[box] (tr) at (0,0) {transactions\\2026-09-22};
\node[box,fill=omMeth!25] (cl) at (3,0.9) {clôtures\\(corrigées)};
\node[box] (ba) at (3,-0.9) {barres à la minute};
\node[box,fill=omThm!12] (mk) at (6,1.5) {valorisations};
\node[box,fill=omThm!12] (pn) at (9,1.5) {P\&L};
\node[box,fill=omThm!12] (ft) at (6,-0.3) {variables};
\draw[->] (tr) -- (cl); \draw[->] (tr) -- (ba);
\draw[->] (cl) -- (mk); \draw[->] (mk) -- (pn);
\draw[->] (cl) -- (ft); \draw[->] (ba) -- (ft);
\node[omThm,anchor=west] at (7.2,-0.3) {périmé~: à reconstruire};
\node[anchor=north] at (3,-1.25) {non touchées};
\end{tikzpicture}

\omcaption{Le lignage des résultats du jour. La \omterm{def:pl:data-quality-and-lineage:correction}{correction du fournisseur} sur la clôture rend périmée toute version construite à partir de l'ancienne clôture
(en rouge)~: les valorisations, le P\&L et les variables de recherche~; les barres à la minute, construites à partir des seules transactions, ne sont pas touchées. Données~:
\texttt{pl\_dataqual.correction\_and\_lineage}.}\label{fig:pl:data-quality-and-lineage:lineage}
\end{omfigure}

\begin{dated}{2026-09}{Une norme pour la provenance}\label{dat:pl:data-quality-and-lineage:prov}
Le modèle de données PROV du World Wide Web Consortium (une recommandation du W3C du 30 avril 2013, consultée en septembre 2026) définit la provenance
comme un enregistrement qui décrit les personnes, institutions, entités et activités impliquées dans la production, l'influence ou la livraison
d'un élément de données~; le graphe de lignage de ce chapitre est une petite instance de ses entités (versions de jeux de données) et de ses activités (les travaux
qui les dérivent).
\end{dated}

Ce que vaut la qualité se lit dans les chiffres du jour (graine~0). Les défauts injectés étant conservés, le VWAP de SIM1 est surévalué de 2.69~\% et
celui de SIM2 de 1.08~\%, tandis que la variance réalisée des variations de prix des transactions de SIM2 vaut 24\,000 fois sa valeur réelle~: huit erreurs de décimale pèsent plus
qu'une journée de négociation authentique. Une fois les lignes d'erreur graves exclues, les deux VWAP sont à moins de 0.03~\% de la vérité (le reliquat tient aux transactions
perdues dans le trou, qu'aucun signalement ne peut restaurer) et la variance réalisée est égale à la vraie.

\section{Tutoriel~: injecter, signaler, corriger}\label{tut:pl:data-quality-and-lineage:tut}

\begin{tutorial}
\textbf{Objectif.} Injecter des défauts et des facteurs de confusion dans une journée simulée, exécuter la suite de règles deux fois (naïve et avertie), mesurer la détection et
les faux signalements, et suivre une correction à travers le lignage. \textbf{État final~:}
le \cref{tab:pl:data-quality-and-lineage:detection}, la \cref{fig:pl:data-quality-and-lineage:jump} et les erreurs de VWAP et de variance de la
dernière section.
\begin{tutsteps}
\item \textbf{La journée}~: \texttt{pl\_dataqual.base()} avec sa suspension de cotation et son saut.
\item \textbf{Injection}~: \texttt{plant(seed)} renvoie les tables défectueuses et la vérité sur les lignes qui sont des défauts.
\item \textbf{Lancement}~: \texttt{firm\_dataqual.run(rules(\dots), tables)} (le \cref{lst:pl:data-quality-and-lineage:rules})~; puis
\texttt{score} et \texttt{detection} sur vingt graines, naïve et avertie.
\item \textbf{Effet}~: \texttt{effect(0)}~: VWAP et variance réalisée avec les défauts conservés, avec les lignes signalées exclues, et la vérité.
\item \textbf{Corriger}~: \texttt{correction\_and\_lineage()}~: deux versions du cours de clôture, et les résultats obsolètes.
\end{tutsteps}
\textbf{Pour aller plus loin.} Abaisser le seuil de saut $k$ de 8 à 4 et compter les faux signalements sur la journée naturelle~; injecter un
saut de 2~\% au lieu d'une erreur de décimale, et trouver le plus petit que la règle détecte encore.
\end{tutorial}

\section{Construction~: la couche de qualité des données}\label{bld:pl:data-quality-and-lineage:dataqual}

\begin{build}
\bfield{Rôle} Les contrôles que passe chaque partition de la \omterm{def:pl:a-tick-store-on-open-formats:store}{base de ticks} avant sa publication, les \omterm{def:pl:data-quality-and-lineage:flag}{indicateurs d'anomalie} sur lesquels filtrent les consommateurs, la
\omterm{def:pl:data-quality-and-lineage:flag}{quarantaine} qui retient les partitions défectueuses, et le lignage qui indique ce qu'une correction invalide~; elle s'exécute sous forme de tâches de la
map du chapitre~1 et du scheduler du chapitre~13.
\bfield{Interface} \texttt{Rule(name, kind, severity, owner, params)}, \texttt{Flag(table, row, rule, severity, detail)}~;
\texttt{run(rules\_by\_table, tables)}, \texttt{clean(df, flags, table, severities)}~; \texttt{Quarantine} avec \texttt{apply},
\texttt{release}, \texttt{readable}, \texttt{log}~; \texttt{Lineage} avec \texttt{add}, \texttt{downstream}, \texttt{upstream},
\texttt{stale\_after}~; natures \texttt{sequence}, \texttt{bounds}, \texttt{crossed}, \texttt{spike}, \texttt{stale},
\texttt{cross\_source}, \texttt{busted}.
\bfield{Règles} Les règles sont des données dotées d'un responsable~; les indicateurs ne modifient jamais les lignes~; un indicateur d'erreur met sa partition en \omterm{def:pl:data-quality-and-lineage:flag}{quarantaine}~; une libération nomme une
personne et un motif~; les corrections sont des versions~; les versions de jeux de données sont désignées par le hachage de leur contenu.
\bfield{Tests d'acceptation} \texttt{code/firm/dataqual/tests/}~: chaque nature détecte son défaut injecté et rien d'autre~; la règle de saut confirmée
épargne un décalage de niveau et la règle d'obsolescence une suspension de cotation~; le décompte des \omterm{def:pl:data-quality-and-lineage:cross}{contrôles inter-sources} par tranche~; la \omterm{def:pl:data-quality-and-lineage:flag}{quarantaine} et la libération~; le lignage désigne les
résultats obsolètes d'une correction.
\bfield{Pour aller plus loin} Des règles générées à partir d'un \omterm{def:pl:data-quality-and-lineage:rule}{contrat de données}~; un score de qualité par partition publié avec lui~; des indicateurs stockés sous forme de table de
la \omterm{def:pl:a-tick-store-on-open-formats:store}{base de ticks}, interrogeable comme les données.
\end{build}

\begin{omsources}
\item U.S. Commodity Futures Trading Commission and Securities and Exchange Commission, \emph{Findings regarding the market events
of May~6, 2010}, 30 septembre 2010.
\item W3C, \emph{PROV-DM: The PROV Data Model}, recommandation du W3C, 30 avril 2013.
\end{omsources}

\section{Exercices}

\begin{exercise}[$\star$]\label{exo:pl:data-quality-and-lineage:1}
Les numéros de séquence d'un flux passent de 91\,686 à 91\,691. Combien de messages manquent, et que fait la capture du chapitre~2
à leur sujet~?
\end{exercise}

\begin{exercise}[$\star$]\label{exo:pl:data-quality-and-lineage:2}
Pourquoi la règle de saut a-t-elle un plancher sur son échelle~? Que se passerait-il sans lui sur un instrument à grand pas de cotation dont les transactions
s'exécutent surtout au même prix~?
\end{exercise}

\begin{exercise}[$\star$]\label{exo:pl:data-quality-and-lineage:3}
Une transaction s'exécute à 3\,000\,000 alors que les transactions précédentes étaient à 300\,000 et les suivantes aussi. Quel type de défaut cela est-il probablement
et pourquoi la médiane prospective le confirme-t-elle~?
\end{exercise}

\begin{exercise}[$\star\star$]\label{exo:pl:data-quality-and-lineage:4}
Pourquoi une règle doit-elle signaler plutôt que supprimer~? Citer un consommateur qui a besoin des lignes signalées.
\end{exercise}

\begin{exercise}[$\star\star$]\label{exo:pl:data-quality-and-lineage:5}
Les défauts étant conservés, le VWAP de SIM2 sur la graine~0 est surévalué de 1.08~\%~; après exclusion des lignes d'erreur, il est sous-évalué de 0.02~\%. D'où provient chaque
erreur~?
\end{exercise}

\begin{exercise}[$\star\star$]\label{exo:pl:data-quality-and-lineage:6}
Le cours de clôture officiel est corrigé à 07~h~10 le lendemain matin. Quels résultats de la journée doivent être reconstruits, lesquels n'ont pas à l'être, et comment le lignage
le sait-il~?
\end{exercise}

\begin{exercise}[$\star\star\star$]\label{exo:pl:data-quality-and-lineage:7}
\emph{Programmation.} Lancer \texttt{score} pour la graine~0 avec la suite naïve et la suite avertie. Combien de faux signalements chaque règle déclenche-t-elle, et
quel événement réel est à l'origine de chacun~?
\end{exercise}

\begin{exercise}[$\star\star\star$]\label{exo:pl:data-quality-and-lineage:8}
\emph{Trouver la faille.} «~Notre tâche de contrôle de la qualité des données supprime toute transaction située à plus de 5~\% de la précédente, et tout trou dans les cotations de plus de
30 secondes est comblé en répétant la dernière cotation. Les données nettoyées vont directement dans la \omterm{def:pl:a-tick-store-on-open-formats:store}{base de ticks}.~»
\end{exercise}

\section{Problème~: l'après-midi annulé}

\begin{problem}[{Problème du week-end --- une journée, ses défauts et ses surprises authentiques}]\label{pb:pl:data-quality-and-lineage:1}
La journée simulée du chapitre (une suspension de cotation et un saut de 3~\% de SIM2), ses six défauts injectés et les deux suites de règles.

\medskip\noindent\textbf{Partie I --- Les contrôles.}
\begin{enumerate}
\item Nommer les quatre familles de contrôles et donner une règle de chacune tirée de la suite du chapitre.
\item Qu'est-ce qu'un \omterm{def:pl:data-quality-and-lineage:rule}{contrat de données}, et qu'apporte-t-il aux règles~?
\item Combien la journée compte-t-elle d'événements, de cotations et de transactions~?
\item Lesquels des défauts injectés aucune règle ne peut-elle détecter, à partir des seules données de la journée, et qu'est-ce qui les détecte~?
\item Pourquoi la \omterm{def:pl:data-quality-and-lineage:correction}{correction du fournisseur} sur le cours de clôture n'est-elle pas un \omterm{def:pl:data-quality-and-lineage:flag}{indicateur d'anomalie}~?
\end{enumerate}

\medskip\noindent\textbf{Partie II --- Faux signalements.}
\begin{enumerate}[resume]
\item Pourquoi la règle de saut naïve signale-t-elle des transactions après le saut de 3~\%, et combien sur la graine~0~?
\item Combien de faux signalements de saut sur vingt graines, pour les règles naïve et avertie~?
\item Pourquoi la règle naïve d'obsolescence signale-t-elle la suspension de cotation, et comment la règle avertie l'évite-t-elle~?
\item Qu'aurait fait à la journée du saut un pipeline qui aurait supprimé les lignes signalées~?
\item Pourquoi un test des \omterm{def:pl:data-quality-and-lineage:rule}{règles de qualité des données} doit-il inclure des événements authentiques, et pas seulement des défauts injectés~?
\end{enumerate}

\medskip\noindent\textbf{Partie III --- Ce que vaut la qualité.}
\begin{enumerate}[resume]
\item À quelle distance de la vérité se trouvent les deux VWAP, avec les défauts conservés, et avec les lignes d'erreur exclues~?
\item De quel facteur la variance réalisée de SIM2 est-elle gonflée lorsque les défauts sont conservés~?
\item Que trouve le \omterm{def:pl:data-quality-and-lineage:cross}{contrôle inter-sources} sur vingt graines, et que ne dit-il pas~?
\item En quoi une \omterm{def:pl:data-quality-and-lineage:flag}{quarantaine} diffère-t-elle d'un \omterm{def:pl:data-quality-and-lineage:flag}{indicateur d'anomalie}, et qui la libère~?
\item Quels résultats le cours de clôture corrigé rend-il obsolètes~?
\end{enumerate}

\medskip\noindent\textbf{Partie IV --- Le verdict.}
\begin{enumerate}[resume]
\item Énoncer le \emph{résultat nommé}~: les nombres de détections et de fausses alertes des deux suites sur vingt graines, ainsi que les erreurs de VWAP et de variance
avec et sans indicateurs.
\item Quelle modification unique de la suite naïve a supprimé le plus de fausses alertes~?
\item Qu'aurait contenu la \omterm{def:pl:a-tick-store-on-open-formats:store}{base de ticks} pour l'après-midi du 6 mai 2010, avec et sans la règle sur les transactions défectueuses~?
\item Que doit enregistrer le lignage pour qu'une correction désigne chaque résultat obsolète~?
\item En une phrase~: qu'est-ce qui distingue une couche de qualité des données d'un script de nettoyage~?
\end{enumerate}
\end{problem}

\section{Questions d'entretien}

\begin{interviewq}[$\star$ \iqroles{développeur, chercheur}]\label{iq:pl:data-quality-and-lineage:1}
Quels contrôles effectueriez-vous sur une journée de \omterm{def:pl:capturing-and-storing-tick-data:tick}{données tick} avant que les chercheurs ne les utilisent~?
\end{interviewq}

\begin{interviewq}[$\star\star$ \iqroles{développeur, chercheur}]\label{iq:pl:data-quality-and-lineage:2}
Comment distinguez-vous une mauvaise transaction d'un véritable mouvement de prix~?
\end{interviewq}

\begin{interviewq}[$\star\star$ \iqroles{développeur}]\label{iq:pl:data-quality-and-lineage:3}
Pourquoi signaler les mauvaises données plutôt que de les supprimer~?
\end{interviewq}

\begin{interviewq}[$\star\star$ \iqroles{développeur, risque}]\label{iq:pl:data-quality-and-lineage:4}
Un fournisseur corrige le cours de clôture de la veille. Que doit-il se passer dans les systèmes de la société~?
\end{interviewq}

\begin{interviewq}[$\star\star\star$ \iqroles{développeur}]\label{iq:pl:data-quality-and-lineage:5}
Comment testeriez-vous un système de qualité des données, et comment sauriez-vous qu'il ne signale pas trop d'éléments~?
\end{interviewq}

\begin{interviewq}[$\star\star\star$ \iqroles{développeur}]\label{iq:pl:data-quality-and-lineage:6}
Concevez un lignage pour une plateforme de recherche de sorte que tout résultat publié puisse être remonté à ses données d'entrée et recalculé après une correction.
\end{interviewq}
